\documentclass[11pt,letterpaper]{article}

\usepackage[spanish,es-noshorthands,es-tabla]{babel}
\usepackage{fontspec}
\setmainfont{Source Serif Pro}
% Símbolo del colón, que Source Serif Pro no trae.
\newfontfamily\fuentecolon{DejaVu Sans}[Scale=MatchLowercase]
\newcommand{\colones}{{\fuentecolon ₡}}
\usepackage{microtype}
\usepackage{geometry}
\usepackage{enumitem}
\usepackage{xcolor}
\usepackage{csquotes}
\usepackage{graphicx}
\usepackage{tabularx}
\usepackage{booktabs}
\usepackage{xurl}
\usepackage[colorlinks=true,urlcolor=cEje,linkcolor=black]{hyperref}
\usepackage[breakable]{tcolorbox}

\geometry{margin=2.3cm}
\hbadness=10000
\setlength{\emergencystretch}{1em}
\setlength{\parindent}{0pt}
\setlength{\parskip}{0.5em}
\setlist{nosep}

\definecolor{cEje}{RGB}{26,58,102}
\definecolor{cGris}{RGB}{120,116,110}
\definecolor{cFondo}{RGB}{244,242,238}
\definecolor{cPend}{RGB}{176,48,32}

% Enunciado de cada paso, en gris y con sangría.
\newcommand{\pregunta}[1]{%
  \begin{quote}\color{cGris}\small\itshape #1\end{quote}}

% Marca en rojo lo que falta escribir. Al terminar no debe quedar ninguna
% (make check lo revisa).
\newcommand{\pendiente}[1]{{\color{cPend}\itshape [Pendiente: #1]}}

% Recuadro para un prompt, con fondo suave.
\newenvironment{prompt}{%
  \begin{center}\begin{minipage}{0.92\textwidth}\color{black!80}\small
  \setlength{\parskip}{0.4em}\itshape}{%
  \end{minipage}\end{center}}

% Recuadro con la salida de un prompt, tal como la devolvió la herramienta.
% El texto vive en un archivo aparte, generado con md2tex.py.
%   \resultado{título}{archivo.tex}
\newcommand{\resultado}[2]{%
  \begin{tcolorbox}[breakable,colback=cFondo,colframe=cEje,
    boxrule=0.5pt,arc=1.5pt,left=7pt,right=7pt,top=5pt,bottom=5pt,
    title={#1},fonttitle=\small\bfseries]
  \footnotesize\setlength{\parskip}{0.35em}\raggedright
  \input{#2}
  \end{tcolorbox}}

% Entrada de la biblioteca de prompts.
%   \entrada{nombre}{cuándo usarlo}{plantilla}{buen resultado}{técnicas}
\newcommand{\entrada}[5]{%
  \subsection*{#1}
  \textbf{Cuándo usarlo.} #2

  \textbf{Plantilla.}
  \begin{prompt}#3\end{prompt}

  \textbf{Cómo se ve un buen resultado.} #4

  \textbf{Técnicas.} #5}

\begin{document}

% === Portada =================================================================
\begin{titlepage}
\centering
\vspace*{3.2cm}
{\large\scshape Inteligencia Artificial Generativa\\
Aplicada a la Transformación Digital\par}
\vspace{0.5cm}
{\normalsize Módulo 3. Actividad 3.1\par}
\vspace{2.6cm}
{\color{cEje}\rule{0.55\textwidth}{0.8pt}\par}
\vspace{0.7cm}
{\LARGE\bfseries Mi prompt en la práctica\par}
\vspace{0.5cm}
{\large Seguimiento semanal del mercado de remates judiciales\\
para Subaxa\par}
\vspace{0.7cm}
{\color{cEje}\rule{0.55\textwidth}{0.8pt}\par}
\vfill
{\large Carlos Luis Rojas Aragonés\par}
\vspace{0.3cm}
{\normalsize Chief Technology Officer Program\par}
\vspace{0.3cm}
{\normalsize 10 de noviembre de 2026\par}
\vspace*{1.5cm}
\end{titlepage}

% === Paso 1 (80 a 120 palabras, sin contar los prompts) ======================
\section*{1. El prompt y su iteración}

\pregunta{Describa: (a) la tarea elegida, (b) el prompt final, (c) los
cambios realizados entre la primera y la segunda versión y (d) qué ha
permitido conseguir el resultado mejorado que no permitía el primero.}

% (a) La tarea: qué es, cada cuánto la hace y para qué le sirve.
Cada semana reviso si pasó algo en el mercado de remates judiciales que
afecte a Subaxa: cambios en las fuentes de datos, en los remates o en la
competencia. Ejecuté las dos versiones en Claude Code.

% Primera versión, tal cual la ejecutó. No la corrija ni la traduzca.
\textbf{Primera versión.}
\begin{prompt}
Research consequential developments from the past week in the
judicial-auction market that could affect Subaxa. Prioritize Costa Rica, then
Latin America and international developments with clear transferable
relevance. Focus especially on changes to auction data sources and access,
collection or scraping methods, auction platform practices, and business
models. For each development, explain what changed, why it matters to Subaxa,
and any concrete implication or follow-up worth considering. Use reliable,
current sources and include links. Distinguish confirmed facts from analysis,
avoid routine or speculative news, and say briefly if there were no
consequential developments.
\end{prompt}

% (b) Versión final, íntegra. La rúbrica pide el prompt completo, no un resumen.
\textbf{Versión final.}
\begin{prompt}
You are a market-intelligence analyst reporting to the CTO of Subaxa.

Context: Subaxa is a Costa Rican PropTech company. Its platform helps
investors buy real estate at judicial auctions in Costa Rica: an interactive
map of opportunities, AI-based analysis, and a legal and investment team that
supports the investor through the auction. The platform depends on auction
edicts published in the Boletín Judicial (run by the Poder Judicial) and on
property data from the Registro Nacional, which we collect automatically. I
read this brief every Monday to decide whether we need to change our data
collection, our product or our priorities.

Task: Research developments from November 3 to November 10, 2026 in the
judicial-auction market that could affect Subaxa. Cover Costa Rica first,
then Latin America and other countries only when the lesson clearly transfers
to Costa Rica.

A development is consequential only if it changes at least one of these:
(1) what auction data is published, where, or how it can be accessed (login,
format, fees, terms of use, rate limits); (2) how judicial auctions are run
(electronic auctions, procedural rules, bills in the Asamblea Legislativa
such as 19.722 and 25.745); (3) what competing auction-listing platforms in
Costa Rica offer (launches, pricing, features).

Verification: Do not rely on search-result summaries. Open the primary
sources yourself: the Boletín Judicial site, the Registro Nacional, the Poder
Judicial news site, the Asamblea Legislativa record of those two bills, and
the competitors' sites. Open every article you cite and confirm its date and
what it says. If you could not open a source, say so.

Format: Answer in English, in at most 300 words, not counting links. First
line: the verdict in one sentence (whether there were consequential
developments this week). Then at most three developments, most important
first, each with three labeled lines: \enquote{What changed} (confirmed fact,
with link), \enquote{Why it matters} (your analysis), \enquote{Follow-up}
(one concrete action). \enquote{Watch list}: at most two items from
outside the week, one line each. \enquote{Unverified}: anything you could
not confirm. \enquote{Sources checked}: one line per primary source, saying
whether you opened it and whether you saw any change.

Constraints: Keep confirmed facts separate from your analysis. Leave out
routine or speculative news and anything outside the dates that does not fit
in \enquote{Watch list}. Do not pad: if nothing qualifies, say so and stop
after the fixed sections.
\end{prompt}

% (c) Qué cambió. Una fila por componente; ponga "no estaba" o "igual" donde
% corresponda. Borre las filas que no aporten.
\begin{table}[htbp]\centering\small
\caption{Los cinco componentes en cada versión del prompt.}
\begin{tabularx}{0.95\textwidth}{@{}lXX@{}}
\toprule
Componente & Primera versión & Versión final \\
\midrule
Rol           & No estaba. & Analista de mercado que le reporta al CTO. \\
Tarea         & Novedades de \enquote{la última semana}. & Igual, con fechas exactas. \\
Contexto      & Solo el nombre de Subaxa. & Qué hace Subaxa, de qué fuentes depende y para qué uso el informe. \\
Formato       & Qué explicar de cada novedad. & Máximo 300 palabras, veredicto primero y secciones fijas. \\
Restricciones & Separar hechos de análisis y evitar lo rutinario. & Además, qué cuenta como relevante, abrir las fuentes y no rellenar. \\
\bottomrule
\end{tabularx}
\end{table}

% (d) Qué logró la segunda versión que la primera no. Algo que se pueda ver
% en los dos resultados, no una impresión general.
A la primera versión le faltaban el rol, el contexto y el formato. Su
resultado salió de resúmenes del buscador y se llenó de noticias
viejas. En la segunda definí qué es relevante, pedí abrir las fuentes y
fijé el formato. El nuevo resultado dice cuáles fuentes revisó y cuáles no
pudo leer, y encontró dos cosas que el primero no vio: una circular que fija
el formato de los edictos y los precios de un competidor.

\resultado{Resultado de la primera versión}{resultado-v1.tex}
\resultado{Resultado de la versión final}{resultado-v2.tex}

\clearpage
% === Paso 2 (no cuenta para el límite de palabras) ===========================
\section*{2. Mi biblioteca de prompts}

\pregunta{Cree al menos tres entradas completas que abarquen un mínimo de
dos funciones profesionales o tipos de tareas distintos.}

% Las variables van entre corchetes: [PAÍS], [PERIODO], [AUDIENCIA]...
% Al menos una entrada debe ser de otra función o tipo de tarea.

\entrada{1. Informe semanal de novedades del mercado}
  {Cada lunes, para decidir si una novedad del mercado obliga a cambiar la
  recolección de datos, el producto o las prioridades. Necesita una
  herramienta que pueda buscar en la web y abrir páginas.}
  {You are a market-intelligence analyst reporting to the [PUESTO] of
  [EMPRESA].

  Context: [QUÉ HACE LA EMPRESA Y DE QUÉ FUENTES DE DATOS DEPENDE]. I read
  this brief every [DÍA] to decide [DECISIÓN QUE APOYA EL INFORME].

  Task: Research developments from [FECHA INICIAL] to [FECHA FINAL] in
  [MERCADO] that could affect [EMPRESA]. Cover [PAÍS PRINCIPAL] first, then
  [OTRAS REGIONES] only when the lesson clearly transfers.

  A development is consequential only if it changes at least one of these:
  [CRITERIOS DE RELEVANCIA, NUMERADOS].

  Verification: Do not rely on search-result summaries. Open the primary
  sources yourself: [FUENTES PRIMARIAS]. Open every article you cite and
  confirm its date and what it says. If you could not open a source, say so.

  Format: Answer in [IDIOMA], in at most [N] words. First line: the verdict
  in one sentence. Then at most [N] developments, each with three labeled
  lines: what changed (confirmed fact, with link), why it matters (your
  analysis) and one follow-up action. Then \enquote{Watch list} (at most
  [N] items outside the dates), \enquote{Unverified} and \enquote{Sources
  checked} (one line per primary source).

  Constraints: Keep confirmed facts separate from your analysis. Leave out
  routine or speculative news. Do not pad: if nothing qualifies, say so and
  stop after the fixed sections.}
  {Empieza con el veredicto. Cada hecho trae un enlace que se abrió, no un
  resumen del buscador. Dice qué fuentes revisó y qué no pudo confirmar. Si
  no pasó nada, lo dice en pocas líneas y no rellena.}
  {\emph{Role prompting} (analista de mercado) y \emph{zero-shot}: no lleva
  ejemplos, así que el formato de salida fijo y la definición explícita de
  qué es relevante hacen el trabajo que haría un ejemplo. Usa los cinco
  componentes: rol, tarea, contexto, formato y restricciones.}

\entrada{2. Borrador de artículo para el blog}
  {Cuando quiero explicar un tema técnico a partir de algo que estamos
  haciendo en la empresa, sin revelar información comercial.}
  {You are a technical writer helping the [PUESTO] of [EMPRESA] write for
  [AUDIENCIA].

  Task: Write a blog post about [TEMA]. First explain the concept briefly,
  then show how it works in practice using this experience: [EXPERIENCIA
  REAL, A ALTO NIVEL].

  Context: [QUÉ HACE LA EMPRESA]. We cannot disclose [LO QUE ES
  CONFIDENCIAL]; describe our work only at a high level.

  Format: [N] words, in [IDIOMA], with [N] figures or charts. Include
  [FIGURA ESPECÍFICA, SI HACE FALTA].

  Example: This is a previous post of mine. Match its tone, length and
  structure, not its content: [ARTÍCULO ANTERIOR].

  Constraints: Give a source with a link for every factual or historical
  claim. Mark every number that is illustrative. Do not claim that something
  is new or unique unless a source supports it. End with a list of the
  claims I should verify before publishing.}
  {Se lee de corrido, la audiencia lo entiende y suena como el artículo de
  ejemplo. Cada dato tiene fuente. Las
  cifras de ejemplo están marcadas como ilustrativas y cuadran entre el
  texto y las figuras. No revela nada confidencial.}
  {\emph{Role prompting} (redactor técnico) y \emph{one-shot}: un artículo
  anterior sirve de ejemplo de tono y estructura. Con varios artículos
  pasaría a ser \emph{few-shot}. Las restricciones cubren las fuentes y la
  confidencialidad.}

\entrada{3. Revisión de un texto generado por IA antes de publicarlo}
  {Antes de publicar o enviar cualquier texto hecho con IA que vaya con mi
  nombre o que llegue a inversionistas.}
  {You are a skeptical editor reviewing a text before publication. You did
  not write it and you have no reason to defend it.

  Context: The text below was drafted with an AI tool for [AUDIENCIA]. It
  will be published at [DÓNDE SE PUBLICA] under the name of [AUTOR].
  [EMPRESA] cannot disclose [LO QUE ES CONFIDENCIAL].

  Task: Review it step by step with these four questions: (1) Is it
  accurate? (2) Is it current? (3) Does it fit this context and audience?
  (4) What still has to be verified before acting on it?

  Format: A table with one row per problem: the quoted sentence, which of
  the four questions it fails, why, and the fix you propose. After the
  table, one line with your verdict: publish, publish with changes, or do
  not publish.

  Constraints: Quote the text exactly. Do not rewrite the whole text. Do
  not comment on style unless it changes the meaning. If you cannot check a
  claim, list it as unverified and do not guess.

  Text: [TEXTO A REVISAR]}
  {Cita frases exactas del texto, no impresiones generales. Encuentra
  afirmaciones que suenan bien pero son falsas o no tienen fuente. Separa
  lo que está mal de lo que solo falta verificar. Termina con un veredicto
  claro.}
  {\emph{Role prompting} (editor escéptico, para que no defienda el texto) y
  \emph{chain of thought}: le pido revisar paso a paso con las cuatro
  preguntas del módulo 2 antes de dar el veredicto. La salida va en una
  tabla.}

% === Paso 3 (150 a 200 palabras) =============================================
\section*{3. Reflexión sobre la iteración}

\pregunta{Indique cuál de los cuatro criterios de evaluación (corrección,
pertinencia, formato y tono, e integridad) ha resultado más difícil de
satisfacer. Explique qué revela esta dificultad y cómo cambiará su manera de
utilizar las herramientas de IA.}

% Cuatro ideas, en este orden:
%   1. qué le mostró iterar que no veía con la primera versión;
%   2. cuál criterio costó más, con un ejemplo concreto del resultado;
%   3. si eso es un problema del prompt o un límite de la herramienta;
%   4. qué va a hacer distinto a partir de ahora.
Iterar me mostró que la primera versión dejaba casi todas las decisiones en
manos de la IA: qué era una semana, qué era relevante, qué hacer con lo
viejo y en qué formato responder. El resultado no era malo, pero era el
informe que la IA decidió hacer, no el que yo necesitaba.

La pertinencia y el formato se resolvieron al definir el criterio y fijar la
estructura. El criterio más difícil fue la integridad.
Lo que más me interesa es saber si cambió algo en las fuentes de datos, y
eso casi nunca sale en las noticias. En la segunda versión la IA abrió las
fuentes, pero no pudo leer las ediciones del Boletín Judicial ni entrar al
sitio de la Asamblea.

Creo que esto no es un problema del prompt, sino un límite de la herramienta:
una IA que busca en la web ve lo que se publica, no lo que cambia en
silencio. A partir de ahora voy a pedir siempre la lista de fuentes revisadas
y de lo no verificado, y los cambios técnicos los voy a vigilar en nuestros
propios sistemas, no con un prompt.

\end{document}
